\begin{tikzpicture}[
  pas/.style={caja, font=\diag, text width=3.9cm, minimum width=4.1cm, minimum height=0.55cm},
  chk/.style={caja, font=\diag, text width=4.3cm, minimum width=4.5cm, minimum height=0.6cm, align=left}]

  % Zonas
  \draw[zona] (0,0) rectangle (4.9,-5.4);
  \draw[zona] (5.25,0) rectangle (10.45,-5.4);
  \draw[zona] (10.8,0) rectangle (16,-5.4);
  \node[zonaLabel] at (0,0) {1. SOPORTES, POR REFERENCIA};
  \node[zonaLabel] at (5.25,0) {2. FECHAS DECLARADAS};
  \node[zonaLabel] at (10.8,0) {3. TRANSACCIÓN AUDITADA};

  % 1. Soportes
  \node[cajaClara, font=\diag, text width=3.9cm, minimum width=4.1cm] (r) at (2.45,-0.75) {UUID, versión positiva y SHA-256 esperado};
  \node[pas] (s1) at (2.45,-1.6) {Selecciona el registro dentro del expediente};
  \node[pas] (s2) at (2.45,-2.5) {Verifica el contexto del cifrado y descifra};
  \node[pas] (s3) at (2.45,-3.4) {Compara la huella y valida la evidencia capturada};
  \node[pas] (s4) at (2.45,-4.3) {Admite PDF o DOCX};
  \foreach \a/\b in {r/s1,s1/s2,s2/s3,s3/s4} \draw[flecha] (\a) -- (\b);
  \node[font=\diag, text=ESCOMGris, text width=4.5cm, align=center] at (2.45,-5.0)
    {una pasada por referencia, aunque sirva a dos papeles};

  % 2. Fechas
  \node[cajaOsc, font=\diag, text width=4.3cm, minimum width=4.5cm] (reloj) at (7.85,-0.8)
    {Instante de captura: reloj del servidor, tras preparar los soportes};
  \node[chk] (f1) at (7.85,-1.85) {Rechaza un acto enteramente posterior a la captura};
  \node[chk] (f2) at (7.85,-2.75) {Rechaza una emisión enteramente posterior a la recepción};
  \node[chk] (f3) at (7.85,-3.65) {Intervalos solapados: se aceptan sin inventar un orden interno};
  \draw[flecha] (reloj) -- (f1);
  \draw[flecha] (f1) -- (f2);
  \draw[flecha] (f2) -- (f3);
  \node[font=\diag, text=ESCOMGris, text width=4.6cm, align=center] at (7.85,-4.85)
    {coherencia de lo declarado; no certifica la fecha jurídica del acto};

  % 3. Transacción
  \node[chk] (t1) at (13.4,-1.15)
    {\textbf{Revalida}: cuenta activa, rol, pertenencia, perfil completo, estado activo, etapa vigente y registros documentales};
  \node[chk] (t2) at (13.4,-2.65)
    {\textbf{Captura}: revisión administrativa y su huella, actor (UUID y correo), instante y soportes admitidos};
  \node[cajaOsc, font=\diag, text width=4.3cm, minimum width=4.5cm] (t3) at (13.4,-3.7)
    {Etapa y evento se confirman juntos};
  \draw[flecha] (t1) -- (t2);
  \draw[flecha] (t2) -- (t3);
  \node[cajaError, font=\diag, text width=4.3cm, minimum width=4.5cm, align=left] (t4) at (13.4,-4.65)
    {Revisión obsoleta: conflicto.\\ Soporte sellado entre preparación y confirmación: nueva validación};

  % Flujo entre zonas
  \draw[flecha] (4.9,-2.65) -- (5.25,-2.65);
  \draw[flecha] (10.45,-2.65) -- (10.8,-2.65);
\end{tikzpicture}
